
Table \ref{tab:asGpioPer04a} shows the input and outputs of the GPIO kernel.
\begin{table}[H]
\caption{GPIO Kernel I/O}
\label{tab:asGpioPer04a}
\centering
\begin{tabularx}{\textwidth}{|l |l |l |X|}
  \hline
  Pin & Direction & Width & Explanation \\
  \hline
  \hline
  clk\_i & in & 1 & System clock \\
  \hline
  rst\_i & in & 1 & System reset. Active high. \\
  \hline
  direction\_i(7:0) & in & 8 & From GPIO-BPI (\asgpioDirRegExt register). Defines the direction of each GPIO. If the bit is 0, the GPIO is an output, if the bit is 1, the GPIO is an input. \\
  \hline
  addr\_i(3:0) & in & 4 & From GPIO-BPI. Address for storing cs\_s. \\
  \hline
  data\_i(7:0) & in & 8 &  From GPIO-BPI (\asgpioDatRegExt register). GPIO data in output direction. The elements of the register will be forwarded to a tristate buffer which is active when direction is output and 'Z' when direction is input. \\
  \hline
  data\_o(7:0) & out & 8 &  From pad to GPIO-BPI (\asgpioDatRegExt register). GPIO data in input direction. If the direction is input, the signals from the pads will go through a tristate buffer and will be forwarded to the register. Else, a '0' will be driven. \\
  \hline
  irq\_o(7:0) & out & 8 &  From kernel to GPIO-BPI. A change of an input will be detected and send to the IRQ register. Every possible GPIO input will be routed to an edge detector (both, rising and falling). Those detected edges will be forwarded to the IRQ structure in the BPI. \\
  \hline
  en\_i & in & 1 &  From GPIO-BPI. Enables a write to a GPIO block. It stores the chip select cs\_o.\\
  \hline
  gpio\_io & bidi & 8 &  To inout pins. GPIO inout. \\
  \hline
  cs\_o & out & 1 &  To output pins. Chip select for outside GPIO block. Stored in an internal register (so, synchronized to the data).  \\
  \hline
\end{tabularx}
\end{table}
